\chapter{Principal Component Analysis and Data}

Imagine you have measured a thousand people and written down, for each one, their
height and their arm span. You now have a cloud of a thousand dots floating in a
two-dimensional plane: one axis is height, the other is arm span. If you actually
plotted those dots, you would not see a shapeless blob. You would see a long,
slanted, cigar-shaped smear, stretched out along the diagonal where ``tall'' meets
``long-armed,'' and thin in the perpendicular direction. The cloud has a
\emph{grain} to it, like wood. There is one direction along which the data really
varies, and another, perpendicular to it, along which it barely varies at all.

\textbf{Principal Component Analysis} -- PCA -- finds that grain. It hands you the
directions in which your data spreads out the most, ranked from most spread to
least, together with a number for ``how much spread'' each one carries.

In the previous chapter, we dissected matrices using the Singular Value
Decomposition: every matrix, regardless of its shape, acts by rotating space,
stretching along orthogonal axes, and rotating again ($\mat A = \mat U\mat\Sigma\mat V\T$).
What makes PCA a theorem of linear algebra rather than a purely statistical recipe
is that those directions of greatest spread are the \emph{eigenvectors of a symmetric
covariance matrix} -- or, equivalently and more stably, the \emph{right singular
vectors} of the centered data matrix. Symmetric matrices, the spectral theorem,
and the SVD all come together directly in PCA.

The approach here is slow, geometric, and visual. We learn to \emph{center} a cloud
of data -- sliding it bodily so its balance point sits at the origin -- then build
the \emph{covariance matrix}, the little symmetric matrix that records how the
cloud spreads and how features vary together, and watch its eigenvectors pick out
the grain. From there: what ``variance explained'' means and how it lets us throw
away dimensions without throwing away meaningful information; the link to the SVD
of the centered data matrix, which is the form a numerical analyst actually uses;
and the modern payoff -- compressing images, building ``eigenfaces,'' and making
sense of the high-dimensional embedding vectors that power modern machine learning.

\chapterline

\section{Data as a cloud of points}

\subsection{The data matrix}

All of data analysis starts with a table. You have $n$ \textbf{samples}
(individual observations, such as people, images, transactions, or days) and for each
one you record $d$ \textbf{features} (the numerical quantities you measured, such as
height, arm span, pixel brightness, or temperature). Stack each sample as a row, and
you get an $n\times d$ \textbf{data matrix}
\[
  \mat{X} \;=\;
  \begin{bmatrix}
    \text{---}\ \vec x_1\T\ \text{---}\\
    \text{---}\ \vec x_2\T\ \text{---}\\
    \vdots\\
    \text{---}\ \vec x_n\T\ \text{---}
  \end{bmatrix}
  \in\R^{n\times d}.
\]
Row $i$ is one sample $\vec x_i\in\R^d$. Column $j$ is one feature measured across
all $n$ samples. Geometrically, each row is a single point living in
$d$-dimensional \emph{feature space} -- the coordinate space where each axis
corresponds to one measured feature -- and the whole dataset is a cloud of $n$
points floating in $\R^d$.

\begin{focusbox}
Keep the two readings of $\mat X$ straight. \textbf{Rows} are points in feature
space $\R^d$ -- this is the geometric picture PCA cares about. \textbf{Columns}
are features, the things you measured. PCA asks a question about the rows (how is
the cloud shaped?) and answers it with a matrix built from the columns (how do
features co-vary?).
\end{focusbox}

\begin{example}\label{ch30:tiny-data}
Here is a tiny dataset we will carry through the chapter: five points in the plane
($n=5$, $d=2$).
\[
  \mat X =
  \begin{bmatrix}
    2 & 1\\
    4 & 3\\
    6 & 5\\
    8 & 5\\
    5 & 1
  \end{bmatrix}.
\]
Each row is a dot at $(x,y)$. Plot them and you see five points trending up and to
the right -- a little cloud with a clear diagonal grain. We will find that grain by
hand before the chapter is done.
\end{example}

\begin{figure}[H]
\centering
\begin{tikzpicture}[scale=0.64,>=Latex]
  % Left Panel: Raw Data
  \begin{scope}[shift={(0,0)}]
    \node[anchor=south, font=\footnotesize\bfseries, spacegray] at (4.5, 6.2) {Raw Data (Shifted Cloud)};
    \draw[->,gray!50] (-0.5,0) -- (9.5,0) node[right, font=\tiny]{$x_1$};
    \draw[->,gray!50] (0,-0.5) -- (0,6.2) node[above, font=\tiny]{$x_2$};
    \foreach \x/\y in {2/1,4/3,6/5,8/5,5/1}{
      \fill[vecblue] (\x,\y) circle (4.5pt);
    }
    % Mean center
    \filldraw[vecred] (5,3) circle (3pt);
    \draw[vecred, thick] (4.6,3) -- (5.4,3);
    \draw[vecred, thick] (5,2.6) -- (5,3.4);
    \node[vecred, font=\footnotesize, fill=white, inner sep=1.5pt, rounded corners=1pt] at (5,3.8) {mean $\vec{\mu}=(5,3)$};
    \node[vecblue, font=\tiny] at (2,1) [below left] {$(2,1)$};
    \node[vecblue, font=\tiny] at (8,5) [above right] {$(8,5)$};
  \end{scope}

  % Centering arrow
  \draw[->, thick, quotientorange] (10.2, 3) -- (11.6, 3)
    node[midway, above, font=\scriptsize] {subtract $\vec{\mu}$}
    node[midway, below, font=\tiny\itshape] {center};

  % Right Panel: Centered Data
  \begin{scope}[shift={(16.8,3.0)}]
    \node[anchor=south, font=\footnotesize\bfseries, spacegray] at (0, 3.6) {Centered Data ($\vec{\mu}=\vec{0}$)};
    \draw[->,gray!50] (-4.5,0) -- (4.5,0) node[right, font=\tiny]{$\tilde{x}_1$};
    \draw[->,gray!50] (0,-3.2) -- (0,3.1) node[above, font=\tiny]{$\tilde{x}_2$};
    % Centered points: (2-5, 1-3)=(-3,-2), (4-5, 3-3)=(-1,0), (6-5, 5-3)=(1,2), (8-5, 5-3)=(3,2), (5-5, 1-3)=(0,-2)
    \foreach \x/\y in {-3/-2,-1/0,1/2,3/2,0/-2}{
      \fill[targetgreen] (\x,\y) circle (4.5pt);
    }
    \filldraw[vecred] (0,0) circle (3pt);
    \draw[vecred, thick] (-0.4,0) -- (0.4,0);
    \draw[vecred, thick] (0,-0.4) -- (0,0.4);
    \node[vecred, font=\footnotesize, fill=white, inner sep=1.5pt, rounded corners=1pt] at (1.1,-0.6) {origin $\vec{0}$};
    \node[targetgreen, font=\tiny] at (-3,-2) [below left] {$(-3,-2)$};
    \node[targetgreen, font=\tiny] at (3,2) [above right] {$(3,2)$};
  \end{scope}
\end{tikzpicture}
\caption{The five-point data cloud of Example~\ref{ch30:tiny-data}. \emph{Left:} raw data points
clustered around sample mean $\vec{\mu}=(5,3)$. \emph{Right:} centering subtracts the mean,
sliding the center of gravity to the origin so that variances and covariances reflect genuine spread
about zero.}
\label{ch30:cloud}
\end{figure}

\subsection{The mean, and why we center}

The first thing PCA does is move the cloud so that its center of mass sits at the
origin. This step is so simple it is easy to dismiss, but it is the difference
between measuring how data \emph{spreads} and measuring how big its numbers happen
to be.

The \textbf{mean} (or centroid) of the cloud is the average of all the points,
\[
  \bar{\vec x} \;=\; \frac{1}{n}\sum_{i=1}^{n}\vec x_i
  \;=\;\frac{1}{n}\,\mat X\T\vec 1,
\]
where $\vec 1\in\R^n$ is the all-ones vector (summing a matrix's rows is exactly
right-multiplying its transpose by $\vec 1$). It is itself a point in $\R^d$ -- the
balance point of the cloud, where it would sit perfectly level on a pin.

\textbf{Centering} the data means subtracting the mean from every point:
\[
  \tilde{\vec x}_i \;=\; \vec x_i - \bar{\vec x},
  \qquad
  \mat B \;=\; \mat X - \vec 1\,\bar{\vec x}\T.
\]
The centered matrix $\mat B$ holds the same cloud, slid bodily over so its centroid
lands on the origin. Distances between points, angles, the whole shape -- all
unchanged. Only the location moved.

\begin{intuitionbox}
PCA is about \emph{shape}, not \emph{position}. A cloud of stars looks the same
whether you describe it from Earth or from Mars; only the coordinates of its center
change. Centering throws away the arbitrary ``where is the origin'' information so
that everything left over is about the cloud's intrinsic spread. Skipping this step
is the single most common way to get nonsense out of PCA.
\end{intuitionbox}

\begin{example}\label{ch30:center-it}
Center the five points of Example~\ref{ch30:tiny-data}. The column means are
\[
  \bar x = \tfrac{2+4+6+8+5}{5} = \tfrac{25}{5} = 5,
  \qquad
  \bar y = \tfrac{1+3+5+5+1}{5} = \tfrac{15}{5} = 3,
\]
so $\bar{\vec x}=(5,3)$. Subtracting it from every row gives
\[
  \mat B =
  \begin{bmatrix}
    2-5 & 1-3\\
    4-5 & 3-3\\
    6-5 & 5-3\\
    8-5 & 5-3\\
    5-5 & 1-3
  \end{bmatrix}
  =
  \begin{bmatrix}
    -3 & -2\\
    -1 & 0\\
    1 & 2\\
    3 & 2\\
    0 & -2
  \end{bmatrix}.
\]
As a check, every column of $\mat B$ now sums to zero: $(-3-1+1+3+0)=0$ and
$(-2+0+2+2-2)=0$. That is exactly what ``centered at the origin'' means.
\end{example}

\begin{trapbox}
A centered column sums to zero -- always check this. If after subtracting the mean
your columns do not sum to (numerically) zero, you subtracted the wrong thing,
subtracted a row mean instead of a column mean, or transposed your matrix. Every
downstream formula in this chapter quietly assumes $\mat B$ is centered.
\end{trapbox}

\section{The covariance matrix}

\subsection{Variance along one feature}

Take a single centered feature -- one column of $\mat B$, say the first, with
entries $b_1,\dots,b_n$ (each already has mean zero). Its \textbf{variance} is the
average squared entry,
\[
  s^2 \;=\; \frac{1}{n}\sum_{i=1}^{n} b_i^2
  \;=\; \frac{1}{n}\,\norm{\text{column}}^2.
\]
Variance is just ``average squared distance from the mean.'' A big variance means
the feature is spread out; a small variance means it clusters tightly. Notice it is
a squared length divided by $n$ -- a dot product of a column with itself.

(A note on the divisor: statisticians often use $\tfrac{1}{n-1}$ instead of
$\tfrac1n$, which makes the variance an unbiased estimate of the true population
variance. It changes every number below by the same constant factor $\tfrac{n}{n-1}$
and so changes \emph{none} of the directions PCA finds. We use $\tfrac1n$ for clean
arithmetic and will flag where the choice matters.)

\subsection{Covariance: do two features move together?}

Two features can each be spread out, but the interesting question is whether they
spread \emph{together}. When height is above its mean, is arm span usually above
its mean too? The number that measures this is the \textbf{covariance} -- the
average product of how far two variables stray from their respective means. For two
centered columns $\vec a$ and $\vec b$ of $\mat B$,
\[
  \Cov(\vec a,\vec b) \;=\; \frac{1}{n}\sum_{i=1}^n a_i b_i
  \;=\; \frac{1}{n}\,\vec a\cdot\vec b.
\]
There is the dot product again. Positive covariance: the features tend to be above
or below their means \emph{together} (the cloud tilts up-and-to-the-right).
Negative covariance: when one is high the other tends to be low (the cloud tilts
down-and-to-the-right). Covariance near zero: no linear tilt; the cloud is round or
axis-aligned. And covariance of a feature with \emph{itself} is exactly its
variance.

\begin{intuitionbox}
Covariance is an inner product normalized by sample size, measuring directional alignment.
Two centered features have large positive covariance for the same
reason two arrows have a large positive dot product -- they point ``the same way''
across the samples. Variance is the special case where a feature is dotted with
itself, the squared length. The entire covariance matrix is the table of all these
pairwise dot products, scaled by $1/n$.
\end{intuitionbox}

\subsection{Packing every covariance into one matrix}

We have $d$ features, so there are $d^2$ pairwise covariances (including each
feature with itself). Arrange them in a $d\times d$ grid: entry $(j,k)$ is the
covariance of feature $j$ with feature $k$. This is the \textbf{covariance matrix}
$\mat C$ -- a square symmetric table holding the variance of each feature along its
main diagonal and the pairwise covariances off the diagonal:
\[
  \boxed{\;\mat C \;=\; \frac{1}{n}\,\mat B\T \mat B\;}
  \qquad\in\R^{d\times d}.
\]
Stare at that. $\mat B\T\mat B$ multiplies every column of $\mat B$ against every
other column; entry $(j,k)$ is (column $j$) $\cdot$ (column $k$), which is exactly
$n$ times the covariance. Divide by $n$ and you have the covariance matrix. The
diagonal entries are the variances of the individual features; the off-diagonal
entries are the covariances between feature pairs.

\begin{formalbox}
For a centered data matrix $\mat B\in\R^{n\times d}$, the \textbf{covariance
matrix} is $\mat C=\tfrac1n\mat B\T\mat B$. It is
\begin{enumerate}
  \item \textbf{symmetric:} $\mat C\T=\tfrac1n(\mat B\T\mat B)\T=\tfrac1n\mat B\T\mat B=\mat C$,
        because $\Cov(\vec a,\vec b)=\Cov(\vec b,\vec a)$;
  \item \textbf{positive semidefinite:} for any direction $\vec w\in\R^d$,
        \[
          \vec w\T\mat C\,\vec w=\tfrac1n\vec w\T\mat B\T\mat B\vec w
          =\tfrac1n\norm{\mat B\vec w}^2\;\ge\;0 .
        \]
\end{enumerate}
\end{formalbox}

These two properties are not technicalities; they are the whole reason PCA works.
\emph{Symmetric} means the spectral theorem applies: $\mat C$ has a full set of
\emph{orthonormal} eigenvectors with \emph{real} eigenvalues. \emph{Positive
semidefinite} means those eigenvalues are all $\ge 0$ -- which is exactly what we
need, because each one is about to be revealed as a variance, and a variance can
never be negative.

\begin{connectionbox}
The quantity $\vec w\T\mat C\vec w$ in the positive-semidefinite check is a
\textbf{quadratic form}, the same object you met with positive definite matrices
and will meet again in optimization. Here it has a direct geometric interpretation: $\vec w\T\mat C\vec
w$ is the variance of the data \emph{after projecting it onto the direction}
$\vec w$ (when $\vec w$ is a unit vector). The covariance matrix is precisely the
machine that turns a direction into ``how much does the cloud spread that way.''
Maximizing that variance over unit vectors is an eigenvalue problem for the symmetric matrix $\mat C$.
\end{connectionbox}

\begin{example}\label{ch30:cov-compute}
Compute the covariance matrix of our centered $\mat B$ from
Example~\ref{ch30:center-it}:
\[
  \mat B =
  \begin{bmatrix}
    -3 & -2\\ -1 & 0\\ 1 & 2\\ 3 & 2\\ 0 & -2
  \end{bmatrix}.
\]
We need $\mat B\T\mat B$, a $2\times2$ matrix. Call the columns $\vec a$ (first)
and $\vec b$ (second).
\[
  \vec a\cdot\vec a = 9+1+1+9+0 = 20,\quad
  \vec b\cdot\vec b = 4+0+4+4+4 = 16,
\]
\[
  \vec a\cdot\vec b = (-3)(-2)+(-1)(0)+(1)(2)+(3)(2)+(0)(-2)=6+0+2+6+0=14.
\]
So
\[
  \mat B\T\mat B=\begin{bmatrix}20 & 14\\ 14 & 16\end{bmatrix},
  \qquad
  \mat C=\frac15\begin{bmatrix}20 & 14\\ 14 & 16\end{bmatrix}
        =\begin{bmatrix}4 & 2.8\\ 2.8 & 3.2\end{bmatrix}.
\]
Read it off: feature 1 has variance $4$, feature 2 has variance $3.2$, and they
covary positively ($2.8>0$) -- the cloud tilts up to the right, exactly as
Figure~\ref{ch30:cloud} showed.
\end{example}

\begin{recoverybox}
\textbf{Pause and think: Scaling the cloud.}
Suppose you multiply all entries in your raw dataset by $10$ (for instance,
converting measurements from centimeters to millimeters).
\begin{itemize}
  \item What happens to the mean $\bar{\vec x}$ and the centered matrix $\mat B$?
  \item What happens to the covariance matrix $\mat C$ and its eigenvalues?
  \item What happens to the principal component directions $\vec v_k$?
\end{itemize}
\emph{Answer:} The sample mean $\bar{\vec x}$ and the centered rows of $\mat B$
are scaled by $10$. The covariance matrix $\mat C = \tfrac1n\mat B\T\mat B$ is
scaled by $10^2 = 100$, so every eigenvalue $\lambda_k$ grows by a factor of $100$
(since variance is a squared distance). But the eigenvectors $\vec v_k$ do not
change at all! The grain of the cloud depends on its geometric shape and
orientation, not on whether you measure in centimeters or millimeters.
\end{recoverybox}

\begin{labbox}
Build covariance intuition with three hand clouds, each centered. (a) Four points
$(\pm1,0),(0,\pm1)$: a round cloud -- compute $\mat C$ and watch the off-diagonal
vanish. (b) Four points $(\pm1,\pm1)$ with matching signs $(1,1),(-1,-1)$ doubled:
strong positive covariance. (c) The same with opposite signs: strong negative
covariance. In each case predict the sign of the off-diagonal entry \emph{before}
you multiply $\mat B\T\mat B$, then check.
\end{labbox}

\section{Principal components: the directions of greatest spread}

\subsection{The variance of a projection}

Here is the central question of the whole chapter, stated geometrically. Pick a
unit direction $\vec w$ (a unit vector in $\R^d$). Project every centered point
onto the line through $\vec w$. The projection of point $\tilde{\vec x}_i$ is the
scalar $\tilde{\vec x}_i\cdot\vec w = \tilde{\vec x}_i\T\vec w$ -- its signed
coordinate along that line. We get $n$ numbers, one per sample. \emph{How spread
out are they?}

Because the data is centered, those projected coordinates also have mean zero (the
average of $\tilde{\vec x}_i\T\vec w$ is $\bar{\tilde{\vec x}}\T\vec w=\vec 0\T\vec
w=0$). So their variance is just the average of their squares:
\[
  \Var(\text{projection onto }\vec w)
  =\frac1n\sum_{i=1}^n(\tilde{\vec x}_i\T\vec w)^2
  =\frac1n\norm{\mat B\vec w}^2
  =\vec w\T\Big(\tfrac1n\mat B\T\mat B\Big)\vec w
  =\vec w\T\mat C\,\vec w .
\]
That short calculation is the hinge of PCA, so let us say in words what it says in
symbols: \textbf{the variance of the data along a unit direction $\vec w$ equals
$\vec w\T\mat C\vec w$.} The covariance matrix, sandwiched between a direction and
itself, returns the spread in that direction.

\begin{figure}[H]
\centering
\begin{tikzpicture}[scale=0.55,>=Latex]
  \draw[->,gray!40] (-4.5,0) -- (4.5,0);
  \draw[->,gray!40] (0,-3.5) -- (0,3.5);
  % the line through w
  \draw[structurepurple,very thick] (-3.8,-2.53) -- (3.8,2.53);
  \node[structurepurple] at (3.6,2.9) {direction $\vec w$};
  % centered points
  \foreach \x/\y in {-3/-2,-1/0,1/2,3/2,0/-2}{
    \fill[vecblue] (\x,\y) circle (4.5pt);
    % foot of projection onto line dir (3,2)/sqrt(13)
    \pgfmathsetmacro{\t}{(\x*3+\y*2)/13}
    \pgfmathsetmacro{\px}{3*\t}
    \pgfmathsetmacro{\py}{2*\t}
    \draw[gray!60,dashed] (\x,\y) -- (\px,\py);
    \fill[vecred] (\px,\py) circle (3pt);
  }
  \node[vecred] at (-2,-2.6) {projections};
\end{tikzpicture}
\caption{Projecting the centered cloud onto a candidate direction $\vec w$. The red
dots are the projected coordinates $\tilde{\vec x}_i\T\vec w$. PCA seeks the
direction that makes those red dots as spread out as possible -- the direction
maximizing $\vec w\T\mat C\vec w$.}
\label{ch30:project}
\end{figure}

\subsection{Maximize the spread: an eigenvalue problem}

We want the direction of \emph{greatest} spread, the grain of the cloud. So we
maximize $\vec w\T\mat C\vec w$ over all unit vectors:
\[
  \max_{\norm{\vec w}=1}\ \vec w\T\mat C\,\vec w .
\]
Without the constraint $\norm{\vec w}=1$ this has no answer -- you could scale
$\vec w$ up forever. The constraint is what makes ``direction'' meaningful: we ask
only about the \emph{way} $\vec w$ points, not how long it is.

Now recall the spectral theorem. Because $\mat C$ is symmetric, it has orthonormal
eigenvectors $\vec v_1,\dots,\vec v_d$ with real eigenvalues
$\lambda_1\ge\lambda_2\ge\cdots\ge\lambda_d\ge0$ (we sort them largest first, and
they are $\ge 0$ by positive semidefiniteness). Write any unit vector in this
eigenbasis as $\vec w=\sum_j c_j\vec v_j$ with $\sum_j c_j^2=1$. Then, using
$\mat C\vec v_j=\lambda_j\vec v_j$ and orthonormality,
\[
  \vec w\T\mat C\,\vec w
  =\Big(\sum_j c_j\vec v_j\Big)\T\Big(\sum_k c_k\lambda_k\vec v_k\Big)
  =\sum_j \lambda_j c_j^2 .
\]
This is a weighted average of the eigenvalues, with weights $c_j^2$ that sum to $1$.
A weighted average is largest when all the weight sits on the biggest value: take
$c_1=1$ and the rest zero, i.e. $\vec w=\vec v_1$. Then $\vec w\T\mat C\vec
w=\lambda_1$. So:

\begin{formalbox}\label{ch30:pca-thm}
\textbf{The principal components are the eigenvectors of the covariance matrix.}
Let $\mat C$ have orthonormal eigenvectors $\vec v_1,\dots,\vec v_d$ with
eigenvalues $\lambda_1\ge\cdots\ge\lambda_d\ge0$. Then:
\begin{itemize}
  \item The direction of greatest variance is $\vec v_1$, and that maximum variance
        is $\lambda_1$.
  \item Among directions \emph{orthogonal} to $\vec v_1$, the one of greatest
        variance is $\vec v_2$, with variance $\lambda_2$; and so on down the list.
\end{itemize}
The eigenvectors $\vec v_k$ are the \textbf{principal components} (or principal
directions); the eigenvalues $\lambda_k$ are the \textbf{variances explained} by
each component.
\end{formalbox}

The ``and so on'' is worth spelling out. Why is the next direction $\vec v_2$?
If we require a unit vector $\vec w$ to be orthogonal to the first principal
direction $\vec v_1$, its coordinate along $\vec v_1$ must vanish:
$c_1 = \vec w\cdot\vec v_1 = 0$. The weighted sum of eigenvalues then starts at
the second term:
\[
  \vec w\T\mat C\,\vec w \;=\; \sum_{j=2}^d \lambda_j c_j^2
  \;\le\; \lambda_2 \sum_{j=2}^d c_j^2 \;=\; \lambda_2 ,
\]
with equality achieved when $c_2 = 1$ and all other weights are zero, which means
$\vec w = \vec v_2$. Repeating this orthogonal constraint step by step picks out
each eigenvector down the sorted list. The principal components thus form a complete
\emph{orthonormal} set of axes, ranked by how much of the data's spread each one
carries. They are a rotated coordinate system, custom fit to your cloud.

\begin{intuitionbox}
The covariance matrix stretches space. Feed it the unit circle of directions and it
reports back a variance for each. The largest variance occurs along the top
eigenvector of $\mat C$, which is the long axis of the elliptical cloud. The eigenvectors are the axes of the data
ellipse; the eigenvalues are (the squares of) how far the ellipse reaches along
each axis. PCA is ``find the axes of the ellipse your data traces out.''
\end{intuitionbox}

\begin{figure}[H]
\centering
\begin{tikzpicture}[scale=0.6,>=Latex]
  \draw[->,gray!40] (-4.5,0) -- (4.5,0);
  \draw[->,gray!40] (0,-3.5) -- (0,3.5);
  % data ellipse (rotated)
  \begin{scope}[rotate=34]
    \draw[vecblue,thick] (0,0) ellipse (4 and 1.6);
  \end{scope}
  % principal axes
  \draw[->,vecred,very thick] (0,0) -- (3.31,2.24) node[above right]{$\vec v_1$};
  \draw[->,targetgreen,very thick] (0,0) -- (-0.9,1.33) node[above left]{$\vec v_2$};
  \node[vecred] at (1.9,0.5) {\small$\sqrt{\lambda_1}$};
  \node[targetgreen] at (-1.0,0.45) {\small$\sqrt{\lambda_2}$};
\end{tikzpicture}
\caption{The data ellipse. Its long axis is the first principal component $\vec
v_1$, its short axis the second $\vec v_2$, and the axes are orthogonal because
$\mat C$ is symmetric. The reach along each axis is $\sqrt{\lambda_k}$, the
standard deviation in that direction.}
\label{ch30:ellipse}
\end{figure}

\subsection{Doing it by hand on the five points}

Let us finish the running example: find the principal components of
\[
  \mat C=\begin{bmatrix}4 & 2.8\\ 2.8 & 3.2\end{bmatrix}
\]
from Example~\ref{ch30:cov-compute}. Eigenvalues come from
$\det(\mat C-\lambda\mat I)=0$:
\[
  (4-\lambda)(3.2-\lambda)-2.8^2=0
  \;\Longrightarrow\;
  \lambda^2-7.2\lambda+(12.8-7.84)=0
  \;\Longrightarrow\;
  \lambda^2-7.2\lambda+4.96=0 .
\]
The quadratic formula gives
\[
  \lambda=\frac{7.2\pm\sqrt{7.2^2-4(4.96)}}{2}
        =\frac{7.2\pm\sqrt{51.84-19.84}}{2}
        =\frac{7.2\pm\sqrt{32}}{2}
        =3.6\pm 2\sqrt2 .
\]
So $\lambda_1=3.6+2\sqrt2\approx 6.43$ and $\lambda_2=3.6-2\sqrt2\approx 0.77$.
Sanity checks: the eigenvalues sum to $7.2=\tr\mat C=4+3.2$, and their product is
$\lambda_1\lambda_2=3.6^2-8=12.96-8=4.96=\det\mat C$. Both hold.

For the top eigenvector, solve $(\mat C-\lambda_1\mat I)\vec v_1=\vec 0$. Using the
first row, $(4-\lambda_1)v_{1}+2.8\,v_{2}=0$, so
\[
  v_2/v_1=\frac{\lambda_1-4}{2.8}=\frac{2\sqrt2-0.4}{2.8}\approx\frac{2.428}{2.8}
  \approx0.867 .
\]
Thus $\vec v_1\propto(1,\,0.867)$; normalized, $\vec v_1\approx(0.756,\,0.655)$.
That points up and to the right at about $41^\circ$ -- precisely the grain we eyed
in Figure~\ref{ch30:cloud}. The second component is orthogonal to it,
$\vec v_2\approx(-0.655,\,0.756)$, the thin direction across the cloud.

\begin{labbox}
Verify the eigenvector by hand: multiply $\mat C\,\vec v_1$ with
$\vec v_1=(0.756,0.655)$ and check you get back roughly $6.43\,\vec v_1$.
($\mat C\vec v_1\approx(4\cdot0.756+2.8\cdot0.655,\ 2.8\cdot0.756+3.2\cdot0.655)
=(4.86,\,4.21)$, and $6.43\cdot(0.756,0.655)=(4.86,\,4.21)$. It matches.)
\end{labbox}

\begin{trapbox}
Eigenvectors come with a sign ambiguity: $\vec v_1$ and $-\vec v_1$ are the same
principal \emph{direction} and carry the same variance. Software may hand you
either one, so a PCA plot can appear flipped from one run or library to the next.
Nothing is wrong; an axis has no built-in arrowhead. Only the line matters, not
which end you call positive.
\end{trapbox}

\section{Variance explained and dimensionality reduction}

\subsection{The total variance is the trace}

How much spread does the cloud have \emph{in total}? Add up the variances of all
the features -- the diagonal of $\mat C$ -- which is $\tr\mat C$. But the trace is
also the sum of the eigenvalues. So
\[
  \text{total variance}=\tr\mat C=\sum_{j=1}^d \mat C_{jj}=\sum_{k=1}^d\lambda_k .
\]
This is the conservation law of PCA. Rotating to the principal-component axes does
not create or destroy spread; it only \emph{redistributes} it. The $d$ messy,
correlated feature-variances on the diagonal are repackaged into $d$ clean,
independent component-variances $\lambda_k$ -- same total, but now sorted and
uncorrelated.

The \textbf{fraction of variance explained} by component $k$ is therefore
\[
  \frac{\lambda_k}{\lambda_1+\cdots+\lambda_d} .
\]
In our example, component 1 explains $6.43/7.2\approx 89.3\%$ of all the spread,
and component 2 the remaining $0.77/7.2\approx 10.7\%$. Nearly nine tenths of
everything the five points do is captured by a single direction.

\begin{aestheticbox}
This structural property explains why PCA is so effective in practice. Real high-dimensional datasets are
rarely genuinely high-dimensional. A thousand pixels of a face, a hundred
survey questions, the daily prices of five hundred stocks -- the underlying signal
lives on a thin, gently curved, low-dimensional sheet, and most of the ambient
dimensions carry almost no variance. PCA measures exactly how thin that sheet is.
When the top few eigenvalues hold $95\%$ of the total, you have learned that your
$1000$-dimensional data is, for all practical purposes, only a handful of
dimensions deep. That fact -- the \emph{intrinsic dimension} is small -- is what
makes learning from data possible at all.
\end{aestheticbox}

\subsection{Reducing dimensions: 2D down to 1D}

Now the payoff. If one direction holds almost all the variance, we can describe each
point by its coordinate along that single direction and lose almost nothing. This
is \textbf{dimensionality reduction}: replace each $d$-dimensional point by its few
top principal-component coordinates.

Concretely, to project our five 2D points down to 1D, give each point its single
coordinate along $\vec v_1$:
\[
  z_i=\tilde{\vec x}_i\T\vec v_1 .
\]
Using $\vec v_1\approx(0.756,0.655)$ and the centered rows of $\mat B$:
\[
  \begin{array}{ll}
  (-3,-2):& z=-3(.756)-2(.655)=-2.268-1.310=-3.58,\\
  (-1,\ 0):& z=-1(.756)+0=-0.756,\\
  (\ 1,\ 2):& z=\ \ .756+1.310=2.07,\\
  (\ 3,\ 2):& z=2.268+1.310=3.58,\\
  (\ 0,-2):& z=0-1.310=-1.31.
  \end{array}
\]
Five 2D points have become five single numbers. To \emph{reconstruct} an
approximate point, ride that coordinate back out along $\vec v_1$ and add the mean
back:
\[
  \hat{\vec x}_i=\bar{\vec x}+z_i\,\vec v_1 .
\]
Every reconstructed point lands exactly on the principal line through the centroid;
the part we threw away was its small wobble in the $\vec v_2$ direction. That
discarded wobble has total size $\lambda_2$ -- the $10.7\%$ we chose to forget.

\begin{figure}[H]
\centering
\begin{tikzpicture}[scale=0.62,>=Latex]
  \draw[->,gray!40] (-0.5,0) -- (10,0) node[right]{feature 1};
  \draw[->,gray!40] (0,-0.5) -- (0,7) node[above]{feature 2};
  % principal line through mean (5,3) with direction (0.756,0.655)
  \draw[structurepurple,very thick] (1.54,0) -- (8.5,6.03);
  \node[structurepurple] at (8.4,6.7) {\small PC1 line};
  \foreach \x/\y in {2/1,4/3,6/5,8/5,5/1}{
    \fill[vecblue] (\x,\y) circle (4pt);
    % project onto line: mean (5,3), dir d=(0.756,0.655)
    \pgfmathsetmacro{\t}{(\x-5)*0.756+(\y-3)*0.655}
    \pgfmathsetmacro{\px}{5+\t*0.756}
    \pgfmathsetmacro{\py}{3+\t*0.655}
    \draw[gray!55,dashed] (\x,\y) -- (\px,\py);
    \fill[vecred] (\px,\py) circle (3pt);
  }
  \fill[black] (5,3) circle (2.5pt) node[below right]{\small mean};
\end{tikzpicture}
\caption{Dimensionality reduction from 2D to 1D. Each blue point collapses (dashed)
onto its red reconstruction on the PC1 line. The reconstruction error is the dashed
length, and PCA chooses the line that makes the \emph{total squared} dashed length
as small as possible.}
\label{ch30:reduce}
\end{figure}

\subsection{Two faces of the same optimum}

We motivated principal components as the directions of \emph{maximum variance}. But
Figure~\ref{ch30:reduce} suggests a second story: PCA is the line (or plane) that
sits \emph{closest} to the data, minimizing the total squared distance from points
to their projections. These are the same optimum, and seeing why ties the chapter
to the projection and least-squares chapters.

\begin{formalbox}
For each centered point, Pythagoras on the right triangle ``point, its projection,
the origin'' reads
\[
  \underbrace{\norm{\tilde{\vec x}_i}^2}_{\text{fixed}}
  =\underbrace{(\tilde{\vec x}_i\T\vec w)^2}_{\text{variance captured}}
  +\underbrace{\norm{\tilde{\vec x}_i-(\tilde{\vec x}_i\T\vec w)\vec w}^2}_{\text{squared error}} .
\]
Summing over all points and dividing by $n$: $(\text{total variance})=(\text{variance
along }\vec w)+(\text{mean squared error})$. The left side does not depend on
$\vec w$. So \textbf{maximizing the variance captured is identical to minimizing the
reconstruction error.}
\end{formalbox}

\begin{figure}[H]
\centering
\begin{tikzpicture}[scale=0.88,>=Latex]
  \draw[structurepurple, thick] (-1.35, -0.66) -- (6.8, 3.32) node[above right, font=\footnotesize, structurepurple] {line $\operatorname{span}(\vec w)$};
  \draw[spacegray, thick] (3.640, 1.776) -- (3.487, 2.091) -- (3.802, 2.244);
  \draw[->, targetgreen, very thick] (0,0) -- (3.955, 1.929);
  \node[targetgreen, font=\footnotesize, anchor=north west] at (2.0, 0.8) {projection $(\vec x_i\cdot\vec w)\vec w$};
  \draw[->, trapred, very thick, dashed] (3.955, 1.929) -- (2.903, 4.086);
  \node[trapred, font=\footnotesize, align=left, anchor=west] at (3.5, 3.2) {error\\$\vec x_i - (\vec x_i\cdot\vec w)\vec w$};
  \draw[->, vecblue, very thick] (0,0) -- (2.903, 4.086);
  \node[vecblue, font=\footnotesize, anchor=south east] at (1.35, 2.1) {$\vec x_i$ (fixed length $\norm{\vec x_i}$)};
  \fill[vecblue] (2.903, 4.086) circle (3.5pt) node[above, font=\footnotesize] {point $\vec x_i$};
  \fill[targetgreen] (3.955, 1.929) circle (3pt) node[below right, font=\footnotesize] {projection foot};
  \fill[black] (0,0) circle (2.5pt) node[below left, font=\footnotesize] {origin $\vec 0$};
  \draw[->, structurepurple, line width=1.3pt] (0,0) -- ({1.5*cos(26)}, {1.5*sin(26)}) node[above left, font=\scriptsize] {$\vec w$};
\end{tikzpicture}
\caption{The Pythagorean decomposition for a single centered point $\vec x_i$.
The distance to the origin, $\norm{\vec x_i}$, is fixed. The projection along
unit vector $\vec w$ forms the base of a right triangle, and the reconstruction error
forms the perpendicular leg. Because $\text{base}^2 + \text{height}^2 = \text{hypotenuse}^2$,
tilting the candidate line $\vec w$ to maximize projected spread (variance) is identical
to minimizing perpendicular reconstruction error.}
\label{ch30:pythagoras}
\end{figure}

\begin{strangbox}
Strang teaches least squares as ``project $\vec b$ onto the column space of $\mat
A$'' -- find the point in a \emph{given} subspace closest to your data. PCA is the
striking inverse problem: the subspace itself is unknown, and you choose it to fit
the data best. Least squares asks ``how close is the data to this line?''; PCA asks
``which line is the data closest to?'' Same orthogonality, same Pythagoras, but now
the subspace is the answer instead of the question. Reading this after Strang, you
will recognize the right triangle of the formal box as the very picture from his
projection chapter -- only here we get to rotate the line until the error shrinks to
its minimum.
\end{strangbox}

\subsection{How many components? The scree plot}

To choose how many components to keep, plot the eigenvalues in decreasing order.
This \textbf{scree plot} usually drops steeply and then flattens into a ``rubble''
tail of tiny eigenvalues -- noise directions carrying negligible variance. You keep
the components before the elbow. Equivalently, keep enough components to reach a
target cumulative variance such as $90\%$ or $99\%$.

\begin{figure}[H]
\centering
\begin{tikzpicture}[scale=0.8]
\begin{axis}[
  width=9.5cm,height=5cm,
  ybar, bar width=11pt,
  title={Variance explained by principal component}, title style={font=\small},
  xlabel={component $k$}, ylabel={$\lambda_k$ (variance)},
  ymin=0, xtick=data, xmin=0.3, xmax=8.7,
  axis lines=left, enlarge x limits=0.08,
]
\addplot[fill=vecblue!70,draw=vecblue] coordinates
  {(1,6.4)(2,2.9)(3,1.1)(4,0.45)(5,0.22)(6,0.12)(7,0.07)(8,0.04)};
\end{axis}
\end{tikzpicture}
\caption{A scree plot. The first two or three components tower over a flat tail of
noise. The ``elbow'' after component 3 says: three principal directions describe
this dataset, the rest is rubble.}
\label{ch30:scree}
\end{figure}

\section{PCA is the SVD of the centered data}

\subsection{Why a practitioner never forms $\mat C$}

We built PCA on the covariance matrix $\mat C=\tfrac1n\mat B\T\mat B$. That is
perfect for understanding, but numerically it is a trap: forming $\mat B\T\mat B$
\emph{squares} the data, and squaring squares the spread between large and small
quantities, wrecking precision (a tiny eigenvalue can vanish entirely into rounding
error). The professional route computes the principal components straight from
$\mat B$ using the \textbf{singular value decomposition} -- no covariance matrix
ever formed.

Recall the SVD of the centered data: $\mat B=\mat U\mat\Sigma\mat V\T$, where $\mat
U\in\R^{n\times n}$ and $\mat V\in\R^{d\times d}$ are orthogonal, and $\mat\Sigma$
holds the singular values $\sigma_1\ge\sigma_2\ge\cdots\ge0$ on its diagonal. Now
substitute into the covariance matrix:
\[
  \mat C=\frac1n\mat B\T\mat B
        =\frac1n(\mat U\mat\Sigma\mat V\T)\T(\mat U\mat\Sigma\mat V\T)
        =\frac1n\mat V\mat\Sigma\T\mat U\T\mat U\mat\Sigma\mat V\T
        =\mat V\Big(\frac{\mat\Sigma\T\mat\Sigma}{n}\Big)\mat V\T .
\]
The middle factor $\mat\Sigma\T\mat\Sigma/n$ is diagonal with entries
$\sigma_k^2/n$. But $\mat V(\cdots)\mat V\T$ with a diagonal middle and orthogonal
$\mat V$ is exactly an eigendecomposition. Therefore:

\begin{formalbox}\label{ch30:pca-svd}
Let $\mat B=\mat U\mat\Sigma\mat V\T$ be the SVD of the centered data. Then the
principal components and their variances are
\[
  \boxed{\;\vec v_k = \text{$k$th right singular vector},
  \qquad
  \lambda_k=\frac{\sigma_k^2}{n}.\;}
\]
The \textbf{right singular vectors of the centered data are the principal
components}, and the \textbf{squared singular values (over $n$) are the variances
explained}. The PCA scores -- the coordinates of the data along the components --
are the columns of $\mat U\mat\Sigma$, since $\mat B\mat V=\mat U\mat\Sigma$.
\end{formalbox}

So ``do PCA'' and ``take the SVD of the centered data, read off $\mat V$ and
$\mat\Sigma$'' are the same instruction. Everything from the eigenvalue story
survives, just expressed in singular values -- and computed without ever squaring
the data.

\begin{intuitionbox}
The SVD already told you that every matrix acts by \emph{rotate, stretch, rotate}.
For the centered data $\mat B$, the final rotation $\mat V$ aligns the coordinate
axes with the cloud's own axes, and the stretches $\sigma_k$ say how far the cloud
reaches along each. That is the data ellipse of Figure~\ref{ch30:ellipse} all over
again: the SVD hands you the ellipse's axes ($\mat V$) and radii ($\sigma_k$)
directly, without the detour through covariance. Figure~\ref{ch30:svd-rss} traces
the three motions in order.
\end{intuitionbox}

\begin{figure}[H]
\centering
\begin{tikzpicture}[scale=0.92,>=Latex]
  % Stage 1: unit circle with two marked axes (the V^T frame)
  \begin{scope}[xshift=0cm]
    \draw[gray!40] (0,0) circle (0.9);
    \draw[->,vecblue,thick] (0,0) -- (0.9,0);
    \draw[->,vecred,thick]  (0,0) -- (0,0.9);
    \node[align=center] at (0,-1.5) {\small unit circle\\\small of directions};
  \end{scope}
  \draw[->,gray!70] (1.2,0) -- (2.0,0) node[midway,above]{\small$\mat V\T$};
  % Stage 2: still a circle, but axes rotated (rotate by V^T)
  \begin{scope}[xshift=3.2cm]
    \draw[gray!40] (0,0) circle (0.9);
    \draw[->,vecblue,thick] (0,0) -- ({0.9*cos(35)},{0.9*sin(35)});
    \draw[->,vecred,thick]  (0,0) -- ({0.9*cos(125)},{0.9*sin(125)});
    \node[align=center] at (0,-1.5) {\small rotated\\\small frame};
  \end{scope}
  \draw[->,gray!70] (4.4,0) -- (5.2,0) node[midway,above]{\small$\mat\Sigma$};
  % Stage 3: stretch into an ellipse along those axes
  \begin{scope}[xshift=6.6cm]
    \draw[vecblue,thick,rotate=35] (0,0) ellipse (1.4 and 0.55);
    \draw[->,vecblue,thick] (0,0) -- ({1.4*cos(35)},{1.4*sin(35)});
    \draw[->,vecred,thick]  (0,0) -- ({0.55*cos(125)},{0.55*sin(125)});
    \node[above left, font=\small, inner sep=1pt] at ({0.75*cos(35)},{0.75*sin(35)}) {$\sigma_1$};
    \node[below left, font=\small, inner sep=1pt] at ({0.35*cos(125)},{0.35*sin(125)}) {$\sigma_2$};
    \node[align=center] at (0,-1.5) {\small data\\\small ellipse};
  \end{scope}
  \draw[->,gray!70] (8.1,0) -- (8.9,0) node[midway,above]{\small$\mat U$};
  % Stage 4: final rotation into sample space (schematic re-tilt)
  \begin{scope}[xshift=10.3cm]
    \draw[targetgreen,thick,rotate=-20] (0,0) ellipse (1.4 and 0.55);
    \node[align=center] at (0,-1.5) {\small in sample\\\small space};
  \end{scope}
\end{tikzpicture}
\caption{PCA is the SVD $\mat B=\mat U\mat\Sigma\mat V\T$ read right to left as
\emph{rotate, stretch, rotate}. The rotation $\mat V$ turns the standard axes into
the cloud's own principal axes; $\mat\Sigma$ stretches the unit circle into the data
ellipse, reaching $\sigma_k$ along principal component $\vec v_k$; the final rotation
$\mat U$ places that ellipse among the samples. The principal directions are the
columns of $\mat V$, and the reach of the data ellipse along each is given by the singular values.}
\label{ch30:svd-rss}
\end{figure}

\begin{remark}
The connection runs deeper than convenience. The eigenvectors of $\mat B\T\mat B$
\emph{are} the right singular vectors of $\mat B$, with eigenvalues $\sigma_k^2$ --
a fact you proved when you first built the SVD. PCA is simply that theorem applied
to centered data and renamed in the language of variance. If the SVD chapter felt
abstract, PCA is its concrete reason to exist.
\end{remark}

\subsection{Low-rank approximation is dimensionality reduction}

The SVD view gives the cleanest statement of dimensionality reduction. Truncate the
SVD to its top $r$ terms,
\[
  \mat B \approx \mat B_r=\sum_{k=1}^r \sigma_k\,\vec u_k\vec v_k\T ,
\]
and the Eckart--Young theorem (from the SVD chapter) says $\mat B_r$ is the
\emph{best} rank-$r$ approximation of $\mat B$ in the least-squares sense -- no
other rank-$r$ matrix is closer. Reducing your data to its top $r$ principal
components is exactly forming this best low-rank approximation. The error you incur
is governed by the singular values you dropped:
\[
  \frac{\norm{\mat B-\mat B_r}_F^2}{\norm{\mat B}_F^2}
  =\frac{\sigma_{r+1}^2+\cdots+\sigma_d^2}{\sigma_1^2+\cdots+\sigma_d^2}
  =\frac{\lambda_{r+1}+\cdots+\lambda_d}{\lambda_1+\cdots+\lambda_d}
  =1-(\text{variance explained by top }r).
\]
``Keep $95\%$ of the variance'' and ``make the best rank-$r$ approximation with
$5\%$ squared error'' are literally the same sentence.

\begin{recoverybox}
\textbf{Pause and think: More features than samples ($d \gg n$).}
Suppose you sequence genetic markers for $n=50$ patients, measuring $d=20{,}000$
gene expression levels for each. The covariance matrix $\mat C = \tfrac1n\mat B\T\mat B$
has size $20{,}000\times 20{,}000$. How many nonzero eigenvalues can $\mat C$
possibly have?
\smallskip

\emph{Answer:} At most $n-1 = 49$. Because the $n$ rows of $\mat B$ are centered,
they sum to zero, so $\rank\mat B \le n-1 = 49$. By the rank properties from
Chapter~16, $\rank(\mat B\T\mat B) = \rank\mat B \le 49$. The remaining $19{,}951$
eigenvalues are exactly zero! Computing the SVD of the $50\times 20{,}000$ matrix
$\mat B$ directly finds those $49$ non-trivial directions in a fraction of a second,
whereas forming and diagonalizing a $20{,}000\times 20{,}000$ covariance matrix
would waste gigabytes of memory on thousands of zeros.
\end{recoverybox}

\section{PCA in the wild}

\subsection{Compressing and denoising images}

A grayscale image is a matrix of pixel brightnesses. Take its SVD and keep the top
$r$ singular triplets: you reconstruct a recognizable image from a tiny fraction of
the numbers. An $800\times600$ image has $480{,}000$ pixels, but its rank-$50$
approximation needs only $50\,(800+600+1)\approx 70{,}000$ numbers -- a $7{\times}$
compression -- and often looks nearly identical, because natural images have a few
dominant singular values and a long thin tail. Drop the tail and you have not only
compressed the image but also \emph{denoised} it: random noise spreads its energy
evenly across all singular directions, so the small-$\sigma$ tail you discard is
disproportionately noise.

\begin{labbox}
Make a $3\times3$ image $\mat M$ with rows $(1,2,3),(2,4,6),(3,6,9)$. Notice every
row is a multiple of $(1,2,3)$, so $\rank\mat M=1$: one singular value is nonzero
and the other two are exactly $0$. The rank-$1$ approximation is the image itself,
with zero error -- this picture compresses perfectly. Now nudge one entry by
$0.1$ and watch a small second singular value appear: that is the ``detail'' a
rank-$1$ approximation would round away.
\end{labbox}

\subsection{Eigenfaces}

In the early 1990s, Turk and Pentland built one of the first working face-recognition
systems on nothing but PCA. Treat each face photo as a long vector: an
$N\times N$ image flattened into a point in $\R^{N^2}$ (a $64\times64$ face is a
single point in $4096$-dimensional space). Collect many faces as the rows of a data
matrix, center them by subtracting the \emph{average face}, and run PCA.

The top principal components are themselves images -- reshape each $\vec v_k$ back
into a $64\times64$ grid -- and they are called \textbf{eigenfaces}. The first few
capture the broadest ways faces differ: overall lighting, the shape of the
hairline, the breadth of the jaw, the set of the eyes. Any face is then
approximately the average face plus a short weighted sum of eigenfaces, so a face
that once needed $4096$ numbers is now described by perhaps $50$ coordinates. To
recognize a new face, project it into eigenface coordinates and find the nearest
stored neighbor. Recognition in a $50$-dimensional space, not a $4096$-dimensional
one -- that compression is what made it run on early-1990s hardware.

\begin{applicationbox}
\textbf{PCA inside modern machine learning.} The eigenface concept applies directly
across modern representation learning. Today neural networks map words, images, and user features into
high-dimensional \textbf{embedding} vectors, and those embeddings are routinely
analyzed and compressed with PCA. Three live uses:
\begin{itemize}
  \item \textbf{Visualization.} Embeddings live in hundreds or thousands of
        dimensions. Project onto the top $2$ or $3$ principal components and you can
        \emph{plot} them -- synonyms cluster, categories separate -- a sanity check
        practitioners run constantly. (Nonlinear cousins like t-SNE and UMAP refine
        the picture, but PCA is the plain linear first look.)
  \item \textbf{Whitening and preprocessing.} Rotating data into its
        principal-component axes and rescaling each by $1/\sqrt{\lambda_k}$
        produces \emph{whitened} features with identity covariance -- uncorrelated
        and unit-variance -- which makes downstream training faster and better
        conditioned.
  \item \textbf{Compression at scale.} A retrieval system storing billions of
        embeddings shrinks each one with PCA before indexing, trading a little
        accuracy for a large cut in memory and a faster nearest-neighbor search.
\end{itemize}
The throughline of this book: the linear algebra under the newest models is the
same spectral theorem you can verify on a $2\times2$ matrix by hand.
\end{applicationbox}

\begin{trapbox}
\textbf{Scale your features first.} PCA maximizes numerical variance, and variance
depends on units. If one feature is measured in millimeters and another in
kilometers, the millimeter column has a much larger numerical variance simply
because of the smaller unit, so the first principal component will align almost
entirely with that axis. When features carry different units or wildly different scales,
\emph{standardize} each to unit variance before running PCA (this replaces the
covariance matrix with the \emph{correlation} matrix). Centering removes the
location; standardizing removes the arbitrary choice of units.
\end{trapbox}

\begin{walkbox}
Take a walk and picture a perfectly round cloud of data -- equal variance in every
direction, all eigenvalues equal. What does PCA return? (Every direction is a tie
for ``most spread,'' so the principal components are not unique -- any orthonormal
basis works.) Now squash the cloud slightly along one axis. Suddenly there is a
unique answer. Which way does $\vec v_1$ swing, and why did a tiny squash break a
perfect tie? Sit with how the eigenvectors of a near-multiple of the identity are
exquisitely sensitive -- a first feel for why nearly-equal eigenvalues make PCA
directions unstable, even as the subspace they span stays perfectly stable.
\end{walkbox}

\section{The recipe, start to finish}

\begin{focusbox}
\textbf{PCA in seven steps.}
\begin{enumerate}
  \item Arrange data as $\mat X\in\R^{n\times d}$: rows are samples, columns features.
  \item Center: subtract the column means to get $\mat B$ (and standardize columns
        to unit variance if the features have different units).
  \item Form the covariance $\mat C=\tfrac1n\mat B\T\mat B$ \emph{for understanding};
        in practice, take the SVD $\mat B=\mat U\mat\Sigma\mat V\T$ instead.
  \item Principal components are the eigenvectors of $\mat C$ = right singular
        vectors $\vec v_k$; variances are $\lambda_k=\sigma_k^2/n$, sorted descending.
  \item Choose $r$ from the scree plot or a cumulative-variance target.
  \item Reduce: project each centered point onto the top $r$ components,
        $\vec z_i=\mat V_r\T\tilde{\vec x}_i$.
  \item Reconstruct when needed: $\hat{\vec x}_i=\bar{\vec x}+\mat V_r\vec z_i$, the
        best rank-$r$ fit to the original.
\end{enumerate}
\end{focusbox}

\begin{historybox}
The idea behind principal component analysis was introduced in 1901 by the English
mathematician Karl Pearson in a paper titled \emph{On Lines and Planes of Closest
Fit to Systems of Points in Space}. Pearson was analyzing biometric data---measurements
of living bodies, such as human skull dimensions and limb lengths---and wanted to
summarize dozens of correlated variables without throwing away the essential variation.
Lacking digital computers or modern matrix algebra, Pearson turned to classical
mechanics: he modeled the data points as a cloud of particles of equal mass in physical
space and calculated their principal axes of rotational inertia---the natural directions
about which a spinning physical body turns smoothly without wobbling, exactly as
Euler had analyzed for rigid solids a century earlier. Three decades later, the American
statistician Harold Hotelling (1933) independently developed the algebraic framework of
PCA, defining principal components as the uncorrelated linear combinations of variables
that maximize projected variance.

A major controversy in Pearson's day was the distinction between ordinary least-squares
regression and orthogonal regression. Standard linear regression treats one variable as
special, minimizing vertical distances along an arbitrary dependent axis $y$. Pearson
argued that in observational science and biology, nature does not label one variable as
an ``input'' and another as an ``output''; true geometric fitting requires minimizing
perpendicular distances, treating every feature symmetrically. Another classic pitfall
is the failure to center data: if you forget to subtract the feature means before
computing components, PCA points its first axis directly from the origin toward the
center of the cloud rather than along the true grain of variance. A second classic trap
is scale: because variance depends on units of measurement, measuring one feature in
millimeters and another in kilometers causes the big-number feature to hijack the first
component entirely, requiring features to be standardized to unit variance before
analysis.

What PCA unlocked is optimal linear dimensionality reduction: projecting high-dimensional
data onto a low-dimensional subspace while preserving the maximum possible variance and
decorrelating all features. Yet PCA has honest limitations. It is blind to nonlinear
geometry: if data lies on a curved surface---such as the classic ``Swiss roll,'' a
two-dimensional sheet curled up into a three-dimensional coil---PCA flattens the roll
onto a flat plane, crushing distant points on adjacent layers together. It is also
unsupervised, meaning it inspects only the cloud's geometric spread without knowing any
class labels or categories; as a result, its directions of greatest variance can fail to
separate distinct classes of data. These boundaries spurred the development of nonlinear
manifold learning and neural autoencoders (which we explore in Chapter~34), while the
underlying optimization of variance via the Rayleigh quotient $\vec w\T\mat C\vec w / \norm{\vec w}^2$
takes center stage in Chapter~36.
\end{historybox}

\begin{aestheticbox}
We asked a statistical question -- which direction maximizes variance -- and the
answer is an eigenvector of the sample covariance matrix. We asked for the best
low-rank summary in mean squared error, and the solution is given by the SVD.
Centering, covariance, variance explained, denoising, eigenfaces, and embedding
compression are all practical realizations of the spectral theorem and the
singular value decomposition applied to empirical data matrices.
\end{aestheticbox}

\subsection*{Looking ahead: from data axes to change of basis}

Throughout this chapter, we worked with data points described in the coordinate
system of our measuring tools---centimeters, dollars, pixel values. Those
measurement axes are arbitrary, whereas the cloud has its own principal
directions. Finding $\vec v_1,\dots,\vec v_d$ gives an orthonormal coordinate
frame aligned with the spread of the data.
When we expressed each point $\vec x_i$ by its projected coordinates
$(z_1,\dots,z_d)$ along those principal axes, all
off-diagonal covariances vanished. The variables decoupled completely, and the
covariance matrix in the new basis became purely diagonal.

That is our first live encounter with a foundational theme of linear algebra:
\emph{a change of basis}. In Chapter~31, we step back from empirical data clouds to
general linear transformations. A linear transformation is an intrinsic coordinate-free
mapping between spaces, while a matrix representation depends on the chosen basis.
Just as PCA finds a basis that diagonalizes the empirical covariance matrix, Chapter~31
develops how to construct coordinate systems that make the matrix representation of an
arbitrary linear operator as simple and direct as possible.

\section*{Exercises}
\addcontentsline{toc}{section}{Exercises}

\begin{exercise}\label{ch30:ex-center}
The data matrix is $\mat X=\begin{bmatrix}1&2\\3&6\\5&4\\7&8\end{bmatrix}$.
(a) Find the mean $\bar{\vec x}$. (b) Write the centered matrix $\mat B$ and verify
both columns sum to zero. (c) Compute the covariance matrix $\mat C=\tfrac14\mat
B\T\mat B$.
\end{exercise}

\begin{exercise}\label{ch30:ex-eig}
For the covariance matrix $\mat C=\begin{bmatrix}5&2\\2&2\end{bmatrix}$, find both
eigenvalues and the (normalized) first principal component. What fraction of the
total variance does the first component explain? Check your eigenvalues against
$\tr\mat C$ and $\det\mat C$.
\end{exercise}

\begin{exercise}\label{ch30:ex-project}
Continuing Exercise~\ref{ch30:ex-eig}, suppose a centered point is
$\tilde{\vec x}=\colvec{3\\1}$. (a) Find its coordinate along the first principal
component. (b) Find its rank-$1$ reconstruction $\hat{\vec x}=(\tilde{\vec
x}\T\vec v_1)\vec v_1$. (c) Verify Pythagoras: $\norm{\tilde{\vec
x}}^2=(\tilde{\vec x}\T\vec v_1)^2+\norm{\tilde{\vec x}-\hat{\vec x}}^2$.
\end{exercise}

\begin{exercise}\label{ch30:ex-round}
Show that for a perfectly round cloud the covariance matrix is a multiple of the
identity, $\mat C=s^2\mat I$. Conclude that \emph{every} unit vector is a principal
component with the same variance, so PCA returns no preferred direction. (This is
the walking-question cloud made precise.)
\end{exercise}

\begin{exercise}\label{ch30:ex-trace}
Prove that the total variance of a centered dataset, $\sum_j\Var(\text{feature }j)$,
equals $\sum_k\lambda_k$, the sum of the eigenvalues of $\mat C$. (Hint: both equal
$\tr\mat C$; use that the trace is basis-independent, or that
$\tr\mat C=\tr(\mat V\mat\Lambda\mat V\T)=\tr\mat\Lambda$.)
\end{exercise}

\begin{exercise}\label{ch30:ex-svd}
The centered data has SVD with singular values $\sigma_1=6,\ \sigma_2=3,\
\sigma_3=2$ and $n=10$ samples. (a) What are the three variances $\lambda_k$?
(b) What fraction of variance does a rank-$2$ reduction retain? (c) What is the
relative squared reconstruction error $\norm{\mat B-\mat B_2}_F^2/\norm{\mat
B}_F^2$?
\end{exercise}

\begin{exercise}\label{ch30:ex-maxvar}
Prove the maximum-variance characterization directly. Let $\mat C$ be symmetric
positive semidefinite with top eigenvalue $\lambda_1$ and unit eigenvector
$\vec v_1$. Using the spectral decomposition $\mat C=\sum_k\lambda_k\vec v_k\vec
v_k\T$ and writing a unit vector as $\vec w=\sum_k c_k\vec v_k$, show
$\vec w\T\mat C\vec w=\sum_k\lambda_k c_k^2\le\lambda_1$, with equality iff
$\vec w=\pm\vec v_1$ (assuming $\lambda_1 > \lambda_2$).
\end{exercise}

\begin{exercise}\label{ch30:ex-units}
A dataset has feature 1 in meters (variance $0.4$) and feature 2 in centimeters
(variance $1600$), with covariance $15$. (a) Find the first principal component of
the raw covariance matrix and argue it points essentially along feature 2.
(b) Convert feature 2 to meters and recompute; show the components change
completely. Explain in one sentence why standardizing features first is essential
here.
\end{exercise}

\begin{exercise}\label{ch30:ex-rank1}
Let $\mat B=\vec a\vec w\T$ be a rank-$1$ centered data matrix ($\vec a\in\R^n$
with $\vec 1\T\vec a=0$, and $\vec w\in\R^d$ a unit vector). (a) Show $\mat
C=\tfrac1n\mat B\T\mat B$ has rank $1$. (b) Find its single nonzero eigenvalue and
the corresponding principal component. (c) Interpret: what does rank-$1$ centered
data look like as a cloud of points?
\end{exercise}

\begin{exercise}\label{ch30:ex-recon-line}
Prove that the rank-$1$ PCA reconstructions $\hat{\vec x}_i=\bar{\vec
x}+(\tilde{\vec x}_i\T\vec v_1)\vec v_1$ of any dataset all lie on a single line
through the mean, and that this line is the total-least-squares best-fit line
(minimizing \emph{perpendicular} distances, unlike ordinary regression which
minimizes vertical distances). Contrast this with the least-squares regression line
of an earlier chapter.
\end{exercise}

\begin{exercise}\label{ch30:ex-mental}
\textbf{Mental checks.} Answer each question in a single sentence or quick calculation:
\begin{enumerate}[label=(\alph*)]
  \item If a dataset has $n=200$ samples and $d=4$ features, what are the matrix
        dimensions of $\mat X$, the centered matrix $\mat B$, and the covariance
        matrix $\mat C$?
  \item Can the covariance matrix $\mat C$ ever possess a strictly negative
        eigenvalue? Why or why not?
  \item If the first two principal components explain $78\%$ and $14\%$ of the
        total variance respectively, what percentage of the variance is carried by
        all remaining components together?
  \item Suppose one software package computes the first principal component as
        $\vec v_1$, while a second package outputs $-\vec v_1$ on the same data.
        Did one of them make a calculation error?
\end{enumerate}
\end{exercise}
